\section{Introduction}
\label{sec:intro}

A tenant on a modern SaaS platform issues a database read, an embedding
lookup and a chat completion inside one user-facing transaction, and the
cluster that serves all three is steered by three controllers that have
never met. A replica autoscaler---the Kubernetes Horizontal Pod Autoscaler
\citep{k8shpa}, KEDA \citep{keda} or a research successor
\citep{firm2020,sageserve2025,chiron2025,aladdin2024}---decides how much
capacity each service holds. A semantic response cache
\citep{gptcache2023,scalm2024,instcache2024,meancache2024} decides which AI
answers are served from an embedding-similarity store instead of a model
call. A model-tier router in the INFaaS tradition \citep{infaas2019} decides
which of several differently priced model classes answers each request.
Each layer has a mature literature. In every system we surveyed, each also
has its own controller: an autoscaler that has never heard of the cache, a
cache that does not know what a replica costs, a router that optimises one
request at a time.

The three layers nonetheless share one budget, one latency target and one
set of tenants, and they interact in ways none of the three controllers can
see. Growing the cache absorbs load that would otherwise demand replicas.
Escalating a tenant to a larger model raises per-request cost that a budget
must absorb, and changes the very latency signal the replica controller
scales on. A locally optimal decision at each layer can therefore be
collectively wasteful, and the waste lands on the platform's bill and on
the tenants' SLOs. This matters more every year: inference is now the
dominant line item for many platforms, model-tier prices span two orders of
magnitude, and the shift to agentic workloads---multi-step chains whose
latency is dominated by the model tier rather than by queueing---puts the
layers' interaction on the critical path.

The literature has moved toward joint decisions, but always within a layer.
Aladdin \citep{aladdin2024} co-optimises placement and scaling; horizontal
and vertical scaling have been reconciled \citep{taleoftwoscales2024}; SLO
constrained joint allocation widens a single layer's action space
\citep{jointalloc2026}; INFaaS chooses among model variants. None of these
controls a semantic cache, none carries cross-tenant fairness in its
objective, and none accounts per-tenant spend in dollars against a budget.
The 2025--26 production stack---llm-d \citep{llmd2025}, AIBrix
\citep{aibrix2025}, NVIDIA Dynamo \citep{dynamo2025}, the Gateway API
Inference Extension \citep{gatewayapi2025}---actuates on engine-level
signals (queue depth, KV-cache pressure, prefill/decode split) and leaves
the portfolio decision above it to whoever operates the cluster. Whether
per-layer optima compose into a good joint decision has, as far as we can
find, never been measured.

Our idea is to treat the three knobs as one decision. \polyforge{}'s
planner, \jcac{} (Joint Cache--Autoscaling Control), holds for every tenant
a triple---replicas, semantic-cache budget, model tier---and every ten
seconds solves one receding-horizon problem over a move-blocked action
lattice, minimising a single objective in cost, SLO violation and Jain
fairness under per-tenant budget caps and cluster caps. The solver is two
exact Gauss--Seidel sweeps over tenants; it is deterministic, runs in
microseconds, and is the same code offline and online. A Kubernetes
operator actuates the plan under guardrails.

The idea was tested rather than asserted, and we report what the test
found. Every confirmatory campaign was pre-registered: hypotheses, sole
treatment, analysis script and stopping rule were committed and pushed to a
public repository before the first run---seven protocols at the first
campaign, forty-eight by the last. The baselines were grid-tuned on our own
composite objective. On that objective the joint controller wins in every
campaign that tested it, including against an offline-trained learned
controller over the identical action space, and it is the only
Pareto-undominated system in the matrix. The evaluation also found things
we would rather not have found, and they are here at the same prominence:
two pre-registered attempts to beat tuned reactive scalers on SLO
attainment failed; the comparative cost percentages we first published were
withdrawn once an audit of the baselines found three asymmetries and three
pre-registered re-scorings showed the percentages did not survive; and on a
rented GPU host the published controller lost to its own tier-only
ablation before a corrected controller won.

What separates this work from its neighbours is compact. Against
autoscalers, it controls two knobs they do not have and prices decisions
per tenant in dollars. Against cache papers, the cache budget is a decision
variable rather than an input, and hit \emph{precision} is measured on the
same real prompts as hit rate. Against fairness schedulers, it works at the
control-plane altitude---capacity across tenants---and says so. Against
every surveyed system, it publishes its nulls under protocols frozen before
the data existed.

The paper makes the following contributions, each of which is a claim tied
to the section that carries its evidence:

\begin{itemize}
  \item \textbf{Joint cross-layer control} (Section~\ref{sec:design}). The
        first surveyed controller to co-optimise replicas, semantic-cache
        budget and model tier against one multi-tenant objective, with two
        construction-level guarantees on the deployed solver. Ablating
        jointness worsens the cost term---by $+2884\%$ against the published
        layered comparator, and by a margin 86\% smaller once that
        comparator's absorbing-tier defect is repaired; the smaller number
        is the claim (Section~\ref{sec:eval-matrix}).
  \item \textbf{A measured composition result and its adjudication}
        (Sections~\ref{sec:eval-matrix} and~\ref{sec:eval-adjudication}).
        The composite-objective win against every tuned baseline in every
        campaign; non-dominance in 48 of 60 cells; and the withdrawal of the
        like-for-like cost percentages under three pre-registered
        re-scorings, leaving a feasibility result: the per-tenant budget is
        satisfiable only by a controller holding the tier knob.
  \item \textbf{An honest SLO account} (Sections~\ref{sec:eval-slo}
        and~\ref{sec:eval-forecast}). Two confirmatory failures published as
        failures; violation parity with tuned HPA/KEDA at $-43$ to
        $-48\%$ cost; a confirmed forecasting mechanism whose winner tracks
        measured periodicity across three real traces.
  \item \textbf{Real-data cache results with a quality denominator}
        (Section~\ref{sec:eval-cache}). 29.6\% adaptive hit rate at cosine
        0.85 on LMSYS-Chat-1M, dollar-weighted savings, an empirical
        capacity curve, and---unique among surveyed cache papers---hit
        precision with the stochasticity ceiling that bounds it.
  \item \textbf{Fairness at the control-plane altitude}
        (Section~\ref{sec:eval-fairness}). Wins over a FIRM-style and a
        VTC-style baseline on Jain's index, with the controller's own
        fairness term honestly nulled.
  \item \textbf{A security frontier for shared semantic caches}
        (Section~\ref{sec:eval-security}). A timing side channel at AUC
        0.88 and a defence comparison in which per-tenant partitioning
        dominates padding and TTL jitter, confirmed at chance over the wire.
  \item \textbf{Live evidence at three depths}
        (Section~\ref{sec:eval-live}). A replica-only ordinal check that
        recorded a disagreement, a valid 24-hour soak with two of four
        hypotheses failed, a real trace driving the real cluster, and a
        three-knob live plane on which the published controller lost, the
        corrected controller won at iso-fairness, and the obvious damper
        failed.
  \item \textbf{A pre-registered systems methodology}
        (Section~\ref{sec:method}). Forty-eight protocols pushed before
        their runs, tuned baselines, paired effect sizes with bootstrap
        intervals, substrate separation, adjudication by new record rather
        than by edit, and a published ledger of nulls---a contribution that
        stands independently of any performance number.
\end{itemize}

Section~\ref{sec:threats} confronts the threats to these claims in their
strongest form, and Section~\ref{sec:conclusion} returns to each
contribution with its verdict.
